\DTMsavedate{mydate}{2022-06-30}
\maketitle{Linear Integrated Circuits}{Electronic Engineering}{\DTMusedate{mydate}}

% ----------------------------------------------------- %
% COURSE INFO
\subsection*{Course Information}
\vspace{0ex}%
\noindent\parbox{0.5\textwidth}{%
    \noindent\begin{tabular}{@{}ll}
                 \textsf{Course:}          & 	Linear Integrated Circuits   \\
                 \textsf{Course Code:}         & 	TEE 3111    \\
                 \textsf{Credit Hours:}    & 	10
    \end{tabular}}
\vspace{2ex}\hrule\vspace{2ex}

% ----------------------------------------------------- %
% INSTRUCTOR INFO
\subsection*{Lecturer Information}
\noindent\parbox{0.5\textwidth}{%
    \noindent\begin{tabular}{@{}ll}
                 \textsf{Name:}         &    Svetlana A. Bebova           \\
% \textsf{Phone:}        &    \phone          \\
\textsf{Email:}        &    \href{mailto:swetlana.bebova@nust.ac.zw}{swetlana.bebova@nust.ac.zw} \\
\textsf{Qualifications:}        &  MSc in Radio Electronic Engineering (Technical University of Sofia)
    \end{tabular}}
\vspace{2ex}\hrule\vspace{2ex}


% ----------------------------------------------------- %
% COURSE DESCRIPTION
\subsection*{Course Synopsis}
The module covers operational amplifier circuits: comparators, inverting and non-inverting amplifiers,
mathematical operations, oscillators and multivibrators, active filters; Voltage regulators; Timer ICs and
their applications; Instrumentation amplifiers; Analogue-to-Digital converters and Digital-to-Analogue
converters
\vspace{2ex} \hrule\vspace{2ex}

% ----------------------------------------------------- %
% LEARNING OBJECTIVES
\subsection*{Learning Objectives}
\begin{outline}

    \1 To familiarize students with the most commonly used linear integrated circuits and their configurations, properties, characteristics and applications.

    \1 To introduce the analysis methods of various circuits built on operational amplifiers and linear integrated circuits.

    \1 To introduce the design principles for circuits and systems with linear integrated circuits.

\end{outline}

\vspace{2ex}\hrule\vspace{2ex}

% ----------------------------------------------------- %
% COURSE TOPICS
\subsection*{Topics}
\begin{outline}
    \1 Operational Amplifiers
    \1 Comparators
    \1 Inverting and Non-Inverting Amplifiers
    \1 Mathematical Operations with Op Amps and Linear Integrated Circuits
    \1 Timers and Signal Generators
    \1 Active Filters
    \1 Voltage Regulators
\end{outline}
\vspace{2ex} \hrule\vspace{2ex}

% ----------------------------------------------------- %
% Students Assenssment
\subsection*{Assessment of Students}
\begin{outline}
    \1 The continuous assessment of students is through:
    \begin{enumerate}
        \item two in-class tests, two hours long.
        \item assignments for each topic.
    \end{enumerate}
    \1 There shall be a final examination at the end of the semester.
\end{outline}
\vspace{2ex}\hrule\vspace{2ex}

% ----------------------------------------------------- %
% Grade Evaluation
\subsection*{Course Evaluation and Grading Policies}
Grades are based on:
Continuous assessment (\qty{25}{\percent}),
Final Exam (\qty{75}{\percent})
\vspace{2ex}\hrule\vspace{2ex}

% ----------------------------------------------------- %
% EQUIPMENT
% https://sites.udel.edu/computing-purchases/personal-specs/
% \subsection*{Essential Equipment}

% \vspace{2ex}\hrule\vspace{2ex}

% Policies and Other information
\subsection*{Policies and Other Information}
% Conduct
\subsection*{Key Text}
\begin{itemize}
    \item Integrated Circuits Lecture Notes by  S. A. Bebova, Electronic Engineering Department, NUST
    \item Basic Operational Amplifiers and Linear Integrated Circuits, 1994 by  T.L. Floyd
\end{itemize}
\vspace{2ex}\hrule
\vspace{2ex}\hrule\vspace{2ex}